\begin{claim}{D.2.52}{Proved}{The half-product is not a positive primary operation}
For $L=Y\langle7\rangle$, there is an integral class
$\omega\in H_{15}(L;\ZZ)$ such that
\[
 \langle\theta,\operatorname{red}_2\omega\rangle=1,
 \qquad
 \operatorname{Sq}_*^i(\operatorname{red}_2\omega)=0
                           \quad(i>0).
 \tag{D.2.52a}
\]
In particular the half-product of D.2.45--50 does not belong to
the positive Steenrod-operation image:
\[
 \theta\notin
 \operatorname{Sq}^1H^{14}(L;\FF_2)
 +\operatorname{Sq}^2H^{13}(L;\FF_2)
 +\operatorname{Sq}^4H^{11}(L;\FF_2)
       =\bigl(\mathcal A^{>0}H^*(L;\FF_2)\bigr)^{15}.
 \tag{D.2.52b}
\]
The class $\omega$ is torsion of order at least four.
It is obtained from the top integral cycle of an explicit
finite fifteen-dimensional space mapping to $L$.
This does not assert that $\omega$ is stably spherical:
higher differentials and integral coefficient obstructions
in the stable Hurewicz spectral sequence remain possible.
\end{claim}
\begin{proof}
\textbf{The suspended sphere factor has zero mod-two Hurewicz.}
Put $X_0=\Sigma B$, so $Y=\Sigma X_0$. In the notation
of D.2.22 and D.2.31,
\[
 \pi_{11}Y=\FF_2\{p,s,x,y\},\qquad
 \operatorname{im}\big(E:\pi_{10}X_0\to\pi_{11}Y\big)
                                  \subseteq\langle p,s\rangle.
 \tag{D.2.52c}
\]
Here is the exact argument for the containment. Under once
adjunction to $\Omega Y$, the classes $p,s$ have zero
mod-two Hurewicz image, while the images of $x,y$ are
respectively
\[
              af+fa,\qquad aab+baa
       \quad\text{in }H_{10}(\Omega\Sigma X_0;\FF_2).
       \tag{D.2.52d}
\]
These are the two independent expressions in D.2.31i;
the letters in this display are the suspended James labels.
Their word lengths are two and three. The once-adjoint of
a suspended map $Et$ is the composite of $t$ with the James
unit $X_0\to\Omega\Sigma X_0$. Its homology image has
word length one. Thus neither coefficient of $x,y$ can
occur, proving \textup{(D.2.52c)}.

The connected-cover interface is needed here as well.
For $E_L=\Omega^2L$, the actual calculation D.2.22 gives
\[
                  h_{\ZZ}(p)=4N',\qquad h_{\ZZ}(s)=0
                          \quad\text{in }H_9(E_L;\ZZ),
                  \tag{D.2.52e}
\]
where $N'$ has order eight. Both classes therefore have
zero mod-two Hurewicz already in $E_L$, not merely after
mapping to $\Omega^2Y$. Naturality of the two homology
suspensions proves that their unique lifts to $L$ have
zero ordinary mod-two Hurewicz in degree eleven.
In particular the lift
$\widetilde{Et}:S^{11}\to L$ of the factor in D.2.49a
satisfies
\[
                        h_{\FF_2}(\widetilde{Et})=0.
                        \tag{D.2.52f}
\]
No value for the integral homology suspension of $4N'$
is inferred in this step.

\textbf{An actual finite cone representing a detecting cycle.}
Let $M_4=M(\ZZ/4,7)$, let
$\delta=\partial(2c_4):S^{14}\to M_4$, and let
$\mu:M_4\to L$ be the actual Moore map of D.2.49.
D.2.46 and D.2.49 supply an integer $n$, with orientations
chosen consistently, for which
\[
 \mu\delta=n\bigl(\widetilde\alpha\sigma'
                        +\widetilde{Et}\nu_{11}\bigr)
                               \quad\text{in }\pi_{14}L.
 \tag{D.2.52g}
\]
Only the existence of $n$ is used; its value and the
homotopy class on either side remain undetermined.
Form
\[
\begin{split}
 P_0&=M_4\vee S^7\vee S^{11},\\
 r&=\iota_M\delta-n\iota_7\sigma'-n\iota_{11}\nu_{11},\\
 P&=P_0\cup_r e^{15}.
\end{split}
 \tag{D.2.52h}
\]
The maps $\mu,\widetilde\alpha,\widetilde{Et}$ on the
three summands extend to an actual map $g:P\to L$ by
\textup{(D.2.52g)}. Let $w\in H_{15}(P;\ZZ)=\ZZ$
be the oriented top cellular class, and put $\omega=g_*w$.

We verify its mixed cup coefficient by an actual cofiber
map. Express $r$ as the pinch $S^{14}\to\bigvee_3S^{14}$
followed by the three separate component maps. The induced
cofiber square gives
\[
 P\longrightarrow
 C(\delta)\vee C(-n\sigma')\vee C(-n\nu_{11}),
 \tag{D.2.52i}
\]
with degree one on each top relative cell. The low classes
$g^*a_7$ and $g^*\operatorname{red}_4b_8$ are the pullbacks
of $a_M+a_S$ and $b_M$, respectively. The sphere $S^7$
contributes to $a_7$, but its integral coefficient Bockstein
is zero; the $S^{11}$ summand contributes to neither class.
All cross products between wedge summands in
\textup{(D.2.52i)} vanish.

The Moore-cone coefficient is exactly
\[
       \langle a_M\smile b_M,[C(\delta)]_{15}\rangle=2
                                      \quad\text{in }\ZZ/4.
       \tag{D.2.52j}
\]
This is an actual Stiefel-cone comparison, before composing
the Moore map into the target. The actual
map $c_{2,4}:M(\ZZ/2,7)\to M_4$ has bottom degree two
and top degree one, and
$c_{2,4}l_4=\pm\delta$ by D.2.43.
The source cone is $C(l_4)=V_{9,2}$ by
\cite[paragraph before Lemma 2.7 and Lemma 2.8(ii), p.~523]
{EH5MorisugiMukai2000}. The resulting cone map has top degree one.
The fundamental mod-four bottom class pulls to twice the
mod-two degree-seven class, while the integral Bockstein
reduces to its degree-eight partner. For the sphere bundle
$S^7\to V_{9,2}\to S^8$, the Euler number is two.
Its mod-two Serre sequence therefore gives the two classes
$x_7,y_8$ with $\operatorname{Sq}^1x_7=y_8$ and
$x_7y_8$ the top class. The pulled-back mod-four product
is consequently twice that top class, giving
\textup{(D.2.52j)}.
Equations \textup{(D.2.52i)--(D.2.52j)} now give
\[
 \langle a_7\smile\operatorname{red}_4b_8,
                       \operatorname{red}_4\omega\rangle=2.
 \tag{D.2.52k}
\]
The half-product identity implies the first assertion
of \textup{(D.2.52a)}. D.2.50m then shows that this integral
class has order at least four.

\textbf{Every positive primary square annihilates the class.}
The finite source in \textup{(D.2.52h)} has
\[
 H_{13}(P;\FF_2)=H_{14}(P;\FF_2)=0,
 \qquad H_{11}(P;\FF_2)=\FF_2\{S^{11}\}.
 \tag{D.2.52l}
\]
Thus the first and second homology squares of its top class
are zero. The fourth square lies in the displayed
$S^{11}$ line, whose image under $g$ is zero by
\textup{(D.2.52f)}. The eighth square is zero by instability:
its dual is $\operatorname{Sq}^8$ on degree-seven
cohomology. Larger powers of two cannot contribute.
Naturality therefore annihilates
$\operatorname{red}_2\omega$ under
$\operatorname{Sq}_*^1,\operatorname{Sq}_*^2,
\operatorname{Sq}_*^4,\operatorname{Sq}_*^8$.
Since the Steenrod algebra is generated by the squares
whose indices are powers of two, this proves the full
annihilation in \textup{(D.2.52a)}.

Finally, six-connectivity and instability show that the
entire degree-fifteen positive-operation image is exactly
the sum in \textup{(D.2.52b)}. Every element of that sum
pairs to zero with $\operatorname{red}_2\omega$, whereas
$\theta$ pairs to one. This proves
\textup{(D.2.52b)} without asserting a stable representative
for the top integral cycle of $P$.
\end{proof}
